\section{Introduction and statement}
An interval-closed subset of a finite poset is an order-convex subset: if it contains $x$ and $y$, it contains every $z$ with $x\le z\le y$. A toggle changes membership at one element when the resulting set is still interval-closed, and otherwise fixes the set. To fix the composition convention, let $(x_1,\ldots,x_{mn})$ be a linear extension from minimal to maximal elements and let $t_x$ be the toggle at $x$. We apply $t_{x_{mn}}$ first and $t_{x_1}$ last. Thus, with function composition evaluated right to left,
\[
 R=t_{x_1}\circ t_{x_2}\circ\cdots\circ t_{x_{mn}}.
\] We use the coordinatewise order on $[m]\times[n]$ and write $R$ for this permutation of $\IC([m]\times[n])$.

For a state $I$, let $\per(I)$ be its least positive return time. Set
\[
 A_m(n)=|\IC([m]\times[n])|,\qquad
 E_m(n)=|\{I\in\IC([m]\times[n]):\per(I)\ne n+m+1\}|.
\]
These are numbers of states, not numbers of orbits. Put
\[
 c_m=\frac{\Cat_m}{(2m)!}=\frac1{m!(m+1)!},\qquad
 \Cat_m=\frac1{m+1}\binom{2m}{m}.
\]

\begin{theorem}\label{thm:main}
For every fixed integer $m\ge2$, as $n$ tends to infinity through all positive integers,
\begin{align}
 E_m(n)&=c_m m(m-1)n^{2m-1}+O_m(n^{2m-2}),\label{eq:main-count}\\
 \frac{E_m(n)}{A_m(n)}&=\frac{m(m-1)}n+O_m(n^{-2}).\label{eq:main-ratio}
\end{align}
In particular, the fraction of states of least period $n+m+1$ tends to one.
\end{theorem}

The width is fixed throughout. Neither the implicit constants nor the onset of the asymptotic estimate is claimed uniform in a growing $m$. The theorem does not classify every exceptional orbit. It counts both states that fail to return after $n+m+1$ steps and states whose least period is a proper divisor of that number.

\subsection{Relation to prior work}
Elder, Lafreni\`ere, McNicholas, Striker, and Welch introduced the interval-closed toggle dynamics and established its global description \cite{ELMSW}. Their chain orbit classification, \cite[Theorem 3.2, arXiv v2]{ELMSW}, already implies the width-one density statement recalled in \cref{prop:width-one}. Lafreni\`ere, Lewis, McNicholas, Striker, and Welch simplified that description, classified width two, and proposed the fixed-width density-one statement in Conjecture 4.1 of \cite{LLMSW}. Their Section 4.1 already suggests a separated-endpoint family, Catalan endpoint orders, and comparison with all interval-closed sets. We retain this attribution: the density strategy and the width-two case are prior work.

The row-nonempty enumeration used below is also prior work: \cite[Theorem 2.7, arXiv v1]{ELLMSW} restates the Narayana-number formula from \cite[Theorem 4.4, arXiv v2]{ELMSW}. Noncrossing rotation models belong to the classical rowmotion literature; see \cite{SW}. The contribution developed here is the explicit passage from rotation to interval-closed-set rowmotion when two boundary waves overlap, together with a matching sharp lower bound from actual nonreturn dynamics. The coefficient $2$ at $m=2$ is consistent with, and already recoverable from, the existing width-two classification. No absolute novelty or publication-priority claim follows from the bounded literature comparison accompanying this project.

\subsection{Architecture of the proof}
The proof separates dynamics from enumeration. Set $N=n+m+1$ and $D=3m+8$. We construct injective equivariant maps into the actual state set from:
\begin{itemize}
\item $X$: noncrossing $m$-edge matchings on the $N$-circle with at most one support gap below $D$, excluding edges between ambient-adjacent vertices;
\item $Y$: noncrossing partitions on the $N$-circle with one triple and $m-2$ pairs, all occupied-support gaps at least $D$.
\end{itemize}
Their images are disjoint for all sufficiently large $n$. Almost every source state has least rotation period $N$, and the exceptions have order at most $n^{2m-2}$. Their combined cardinality is
\[
 |X|+|Y|=c_m\bigl(n^{2m}+mn^{2m-1}\bigr)+O_m(n^{2m-2}).
\]
The total state count has second coefficient $m^2$ instead of $m$. Thus these two models give the upper bound in \eqref{eq:main-count}.

For the lower bound, mark a valley $BA$ in a Dyck endpoint word and require the corresponding physical gap to be one. Keep every other endpoint gap and both margins above a fixed width-dependent floor. The resulting states cannot return at time $N$. There are $\frac{m-1}{2}\Cat_m$ marked valleys, and stars and bars gives their total leading coefficient $c_m m(m-1)$. The proof of nonreturn uses a finite sequence of rigorously controlled boundary events; it requires no full classification of their least periods.

The longer local calculations are in the appendices. Their complete candidate phase arrays and cutoff verifications make the dynamical argument independent of finite experiments or the formal certificate.
